\documentclass[11pt]{amsart}

\usepackage[T1]{fontenc}
\usepackage{lmodern}
\usepackage{microtype}
\usepackage{amsmath,amssymb}
\usepackage{graphicx}
\usepackage[hidelinks]{hyperref}

\setcounter{MaxMatrixCols}{12}

\newcommand{\permanent}{\operatorname{per}}
\newcommand{\PhiOne}{\Phi_1}

\newtheorem{theorem}{Theorem}
\newtheorem{lemma}[theorem]{Lemma}
\newtheorem{corollary}[theorem]{Corollary}

\title[A counterexample to the weak Drury inequality]
{A Counterexample to the Weak Drury Permanent Inequality}

\author{Owen Shen}
\address{Operations Research Center, Massachusetts Institute of Technology,
Cambridge, MA 02139, USA}
\email{owenshen@mit.edu}
\date{}

\subjclass[2020]{Primary 15A15; Secondary 15B48, 15A45}
\keywords{permanent, positive semidefinite matrix, Gram matrix, counterexample}

\begin{document}

\begin{abstract}
Zhang recorded as an open problem a weak consequence of the Drury permanent
conjecture for Hermitian positive semidefinite matrices. We give an exact
$12\times12$ Gaussian-integer Gram matrix of rank two for which the proposed
inequality is reversed. Rank two is the least possible rank of a
counterexample, and the least possible order lies between $4$ and $12$.
Positive definite correlation matrix counterexamples exist in every order at
least $12$.
\end{abstract}

\maketitle

\section{The counterexample}

For an $m\times m$ matrix $M=(m_{jk})$, write
\[
  \permanent M=\sum_{\sigma\in S_m}\prod_{j=1}^m m_{j,\sigma(j)}.
\]
For a matrix $A$ of order $n$ and distinct indices $j,k$, let
$A[\widehat{j,k}]$ be the principal submatrix obtained by deleting rows and
columns $j$ and $k$, and set
\begin{equation}\label{eq:phi}
  \PhiOne(A)=\sum_{k=2}^n |a_{1k}|^2\,
  \permanent A[\widehat{1,k}].
\end{equation}
The permanent of the empty matrix is understood to be $1$.

\begin{theorem}\label{thm:main}
Let $i^2=-1$, let $A=V^*V$, and let
\[
\resizebox{0.98\textwidth}{!}{$
V=\begin{pmatrix}
1&1&7&7&9&5&5&7&9&5&5&0\\
0&-2&5&-4+3i&2-6i&8+6i&-3+9i&-4-3i&2+6i&-3-9i&8-6i&1
\end{pmatrix}.$}
\]
Then $A$ is Hermitian positive semidefinite of order $12$ and rank $2$, and
\begin{align*}
  \permanent A
    &=463\,995\,226\,880\,175\,960\,000\,000,\\
  \PhiOne(A)
    &=464\,101\,796\,319\,782\,967\,000\,000.
\end{align*}
Hence
\[
  \PhiOne(A)-\permanent A
  =106\,569\,439\,607\,007\,000\,000>0.
\]
\end{theorem}

Zhang~\cite[p.~315, Eq.~(21)]{Zhang2016} recorded the inequality
\begin{equation}\label{eq:weak}
  \PhiOne(A)\leq\permanent A
\end{equation}
as an open problem for Hermitian positive semidefinite $A$. It is an immediate
consequence of the stronger Drury permanent conjecture
\cite[p.~315, Eq.~(20)]{Zhang2016}, which Hutchinson~\cite{Hutchinson2017}
disproved with a $4\times4$ complex correlation matrix. Following
Hutchinson's indexing, the relevant minor is
$B_{k-1,k-1}=A[\widehat{1,k}]$. Thus Theorem~\ref{thm:main} disproves exactly
Zhang's surviving weak inequality~\eqref{eq:weak}.

As printed, Paksoy~\cite[Eq.~(1.1)]{Paksoy2023} uses order-$(n-1)$ principal
minors in the second term, giving a nonhomogeneous formula different from
both equations~(20) and~(21) of Zhang. The real matrix in
\cite[Eq.~(1.3)]{Paksoy2023} satisfies the intended weak inequality, since
\[
  \PhiOne(A)=\frac{38673}{25000}
  <\frac{2879693}{800000}=\permanent A.
\]

\section{Exact verification}

\begin{lemma}\label{lem:ranktwo}
Let $V=(v_1,\ldots,v_m)\in\mathbb C^{2\times m}$, where
$v_j=(x_j,y_j)^{\mathsf T}$, and put $G=V^*V$. If
\[
  \prod_{j=1}^m(x_j+y_jz)=\sum_{r=0}^m c_rz^r,
\]
then
\begin{equation}\label{eq:ranktwo}
  \permanent G=\sum_{r=0}^m r!(m-r)!\,|c_r|^2.
\end{equation}
\end{lemma}

\begin{proof}
Expand
$\prod_j(\overline{x_j}x_{\sigma(j)}+\overline{y_j}y_{\sigma(j)})$
over subsets $S\subseteq\{1,\ldots,m\}$, where $S$ records the rows in which
the $y$-term is chosen. If $|S|=r$, grouping permutations by
$T=\sigma(S)$ gives
\[
 \sum_{\sigma\in S_m}
 \prod_{j\notin S}x_{\sigma(j)}
 \prod_{j\in S}y_{\sigma(j)}
 =r!(m-r)!\sum_{|T|=r}
 \prod_{k\notin T}x_k\prod_{k\in T}y_k
 =r!(m-r)!c_r.
\]
Summing the conjugate prefactors over all $S$ of size $r$ gives
$\overline{c_r}$, and \eqref{eq:ranktwo} follows.
\end{proof}

\begin{proof}[Proof of Theorem~\ref{thm:main}]
The first and last columns of $V$ are $(1,0)^{\mathsf T}$ and
$(0,1)^{\mathsf T}$, so $A=V^*V\succeq0$ has rank $2$.
Lemma~\ref{lem:ranktwo}, applied to all twelve columns and then after deleting
columns $1$ and $k$, gives the following exact values. Grouped indices have
the same weight and minor permanent.
\[
\begin{array}{c|c|r}
 k & |a_{1k}|^2 & \permanent A[\widehat{1,k}]\\ \hline
 2       & 1  & 9\,076\,601\,601\,430\,046\,100\,000\\
 3       & 49 & 1\,684\,042\,529\,828\,861\,700\,000\\
 4,8     & 49 & 1\,687\,819\,148\,108\,486\,100\,000\\
 5,9     & 81 & 1\,042\,294\,059\,228\,024\,900\,000\\
 6,11    & 25 &   362\,453\,951\,893\,503\,780\,000\\
 7,10    & 25 &   402\,529\,981\,049\,836\,740\,000
\end{array}
\]
Also $a_{1,12}=0$. Substitution into \eqref{eq:phi}, with multiplicity two
for each grouped pair, gives the values in the theorem. Since all input lies
in $\mathbb Z[i]$, the calculation is exact.
\end{proof}

\section{Consequences}

\begin{corollary}\label{cor:minimal}
Rank two is the least possible rank of a counterexample to \eqref{eq:weak}.
If $n_{\min}$ denotes the least possible order, then
$4\leq n_{\min}\leq12$.
\end{corollary}

\begin{proof}
If $A=uu^*$ has rank at most one and
$P=\prod_{j=1}^n|u_j|^2$, then
\[
  \permanent A=n!P,
  \qquad
  \PhiOne(A)=(n-1)!P=\frac1n\permanent A.
\]
For order $2$, $\permanent A-\PhiOne(A)=a_{11}a_{22}\geq0$. For order $3$,
write
\[
 A=\begin{pmatrix}
 a&b&c\\ \overline b&d&e\\ \overline c&\overline e&f
 \end{pmatrix}\succeq0.
\]
Then
\[
 \permanent A-\PhiOne(A)
 =\det A+2a|e|^2+d|c|^2+f|b|^2\geq0.
\]
The theorem supplies the upper bounds.
\end{proof}

\begin{corollary}\label{cor:pd}
For every $N\geq12$, some positive definite Hermitian correlation matrix of
order $N$ violates \eqref{eq:weak}.
\end{corollary}

\begin{proof}
For every positive diagonal $D=\operatorname{diag}(d_1,\ldots,d_n)$,
\[
 \permanent(DAD)=\Bigl(\prod_{j=1}^n d_j^2\Bigr)\permanent A,
 \qquad
 \PhiOne(DAD)=\Bigl(\prod_{j=1}^n d_j^2\Bigr)\PhiOne(A).
\]
The diagonal entries of the matrix in Theorem~\ref{thm:main} are positive, so
diagonal normalization gives a rank-two correlation matrix $A_0$ with the
same strict violation. Set $C=A_0\oplus I_{N-12}$, omitting the identity block
when $N=12$. Then
$\PhiOne(C)-\permanent C=\PhiOne(A_0)-\permanent A_0>0$. The same remains
true for $C+\varepsilon I_N$ when $\varepsilon>0$ is sufficiently small.
This matrix is positive definite, and a final diagonal normalization preserves
the violation.
\end{proof}

\begin{thebibliography}{9}

\bibitem{Hutchinson2017}
G.~Hutchinson,
\emph{A counterexample to the Drury permanent conjecture},
Special Matrices \textbf{5} (2017), 301--302.
\href{https://doi.org/10.1515/spma-2017-0022}
{doi:10.1515/spma-2017-0022}.

\bibitem{Paksoy2023}
V.~E. Paksoy,
\emph{A permanent inequality for positive semidefinite matrices},
Electron. J. Linear Algebra \textbf{39} (2023), 71--77.
\href{https://doi.org/10.13001/ela.2023.7701}
{doi:10.13001/ela.2023.7701}.

\bibitem{Zhang2016}
F.~Zhang,
\emph{An update on a few permanent conjectures},
Special Matrices \textbf{4} (2016), 305--316.
\href{https://doi.org/10.1515/spma-2016-0030}
{doi:10.1515/spma-2016-0030}.

\end{thebibliography}

\end{document}